\chapter{최종 제안: 기대수익이 높은 대출부터 산다}

\begin{keybox}[최종 제안(test 규칙 v2.3, 42단계)]
\textbf{매달 새로 나온 대출마다 기대수익 점수 $s$를 계산하고, 승인비율이 몇 \%든 $s$가 높은 대출부터 산다. 사지 않은 돈은 국채에 둔다.}
\begin{itemize}[leftmargin=*]
\item 쓰는 전략은 \textbf{기대수익 순 하나}다. 몇 \%를 살지는 정하지 않는다. 승인비율은 투자자가 가진 돈에 달려 있고, 우리는 그 안에서 \textbf{무엇을 살지}를 답한다.
\item \textbf{주장:} 승인비율이 10\%든 90\%든, 같은 양을 무작위로 산 것보다 \textbf{더 번다}.
\item \textbf{주 판정:} 처음 규칙서(v1)의 H2를 정의 그대로 쓴다. Sharpe(흔들림)는 보조 지표로 보고한다.
\end{itemize}
\end{keybox}

\begin{formulabox}[승인비율 $q$란]
매달 그 달에 나온 36개월 대출을 점수 순으로 줄 세워 위에서부터 사는 \textbf{건수의 비율}이다(10\%, 20\%, \dots, 90\%). 투자자가 그 달 시장의 몇 \%를 살 돈이 있는지를 나타내며, 사지 않은 돈은 국채에 둔다고 본다. 대출마다 금액이 달라 \textbf{금액 비중} $w$는 $q$와 조금 다르다(예: 부도확률 순 80\%에서 평균 82.5\%). 선별 가치는 승인한 1달러당 IRR 차이로, Sharpe는 금액 비중 $w$를 곱한 초과수익으로 계산한다.
\end{formulabox}

\begin{formulabox}[선별 가치와 포트폴리오 초과수익은 다르다]
\textbf{선별 가치}는 \textbf{대출에 넣은 돈의 질}을 잰다. 어느 달 대출이 1,000건이고 승인비율이 10\%면, 점수 상위 100건과 아무 100건을 각각 사서 대출이 끝난 뒤 두 묶음의 실제 IRR을 비교한다. 예를 들어 우리 묶음 연 8.0\%, 무작위 묶음 연 5.9\%면 선별 가치는 2.1\%p다. IRR은 ``넣은 1달러가 1년에 몇 \% 벌었나''이므로, 얼마를 넣었는지와 상관없이 \textbf{같은 양을 무작위로 산 것보다 1달러당 몇 \%p 더 버나}를 뜻한다. 48개월 평균이 표의 숫자다. 고르는 실력을 재는 숫자이고, 주 판정에 쓴다.

\textbf{포트폴리오 초과수익}은 \textbf{전체 돈의 성과}다. 사지 않은 돈은 국채에 두어 초과수익이 0이므로, 승인비율 10\%면 돈의 90\%가 국채보다 더 벌지 못한다. 투자자가 실제로 국채보다 얼마나 더 버는지는 이 숫자가 보여 준다.

\smallskip
\centering\small
\begin{tabular}{@{}lccccc@{}}\toprule
기대수익 순(V1), train & 10\% & 50\% & 80\% & 90\% & 전부 승인\\\midrule
선별 가치(1달러당, 무작위 대비) & +2.11 & +0.99 & +0.48 & +0.28 & 0\\
포트폴리오 초과수익(국채 대비, 전체 돈) & 0.65 & 2.98 & 4.34 & 4.67 & 4.86\\\bottomrule
\end{tabular}

\raggedright
\smallskip
\textbf{읽는 법:} 적게 살수록 가장 좋은 대출만 골라 1달러당 차이가 크지만, 돈 대부분이 국채에 있어 전체 수익은 작다. 많이 살수록 평범한 대출까지 사게 되어 1달러당 차이는 줄지만 전체 수익은 커진다. 그래서 ``10\%만 사는 것이 가장 좋다''는 뜻이 아니다. 이 기간에는 대출이 국채보다 나았으므로 많이 살수록 전체 수익은 컸고, 우리 주장은 \textbf{몇 \%를 사든 같은 양을 무작위로 사는 것보다 낫다}는 것이다.
\end{formulabox}

\section{흐름: 고르는 날과 채점하는 날}
\begin{figure}[H]\centering
\begin{tikzpicture}[x=1mm,y=1mm,box/.style={draw=black!70,rounded corners=1.2mm,align=center,font=\footnotesize,minimum height=12mm,text width=32mm,fill=white,execute at begin node={\XeTeXlinebreakskip=0pt\relax}},ar/.style={-{Stealth},black!65}]
\node[box] (a) at (0,0) {그 달 새 대출\\(36개월 만기)};
\node[box] (b) at (42,0) {부도확률 $p$\\시점 학습 LightGBM\\(그날 알던 정보만)};
\node[box] (c) at (84,0) {기대수익 점수\\$s=(1-p)r-p\cdot\mathrm{LGD}$\\$-$ 국채 금리};
\node[box,line width=0.9pt] (d) at (126,0) {$s$ 높은 순으로\\예산만큼 구매\\남는 돈은 국채};
\draw[ar] (a)--(b); \draw[ar] (b)--(c); \draw[ar] (c)--(d);
\draw[dashed,black!35] (-18,-13)--(144,-13);
\node[font=\scriptsize\itshape,black!60,anchor=south west] at (-18,-13) {위: 승인하는 날(결과를 모름)};
\node[font=\scriptsize\itshape,black!60,anchor=north west] at (-18,-13) {아래: 대출이 끝난 뒤(실제 결과로 채점, $s$는 쓰지 않음)};
\node[box] (e) at (126,-27) {우리가 산 묶음\\실제 현금흐름\\\ensuremath{\to} 코호트 IRR};
\node[box] (f) at (126,-47) {같은 양 무작위(10회)\\실제 현금흐름\\\ensuremath{\to} 코호트 IRR};
\node[box,line width=0.9pt] (g) at (84,-37) {선별 가치\\$=$ 우리 IRR $-$ 무작위 IRR\\(국채는 양쪽에서 지워짐)};
\node[box] (h) at (42,-37) {48개 발행월 평균\\예산 10--90\%\\블록 부트스트랩 구간};
\node[box,line width=0.9pt] (i) at (0,-27) {주 판정(H2)\\95\% 구간 하한 $>0$\\\ensuremath{\to} 제안 유지};
\node[box,dashed] (j) at (0,-47) {보조 지표\\포트폴리오 Sharpe\\(평균 \ensuremath{\div} 흔들림)};
\draw[ar] (d)--(e); \draw[ar] (e)--(g); \draw[ar] (f)--(g); \draw[ar] (g)--(h); \draw[ar] (h)--(i); \draw[ar,dashed] (h)--(j);
\end{tikzpicture}
\caption{제안의 흐름. 위 줄은 승인하는 날 하는 일, 아래 줄은 대출이 끝난 뒤 채점하는 일이다. 코호트 IRR은 빌려준 돈(0개월), 월 납입금, 마지막 달의 조기상환 일시금이나 부도 후 회수금을 묶음 전체로 더해 구하고, 받은 돈에서 수수료 1\%를 뗀다(2장).}
\end{figure}


\begin{itemize}[leftmargin=*]
\item \textbf{$s$는 줄 세우기에만 쓴다.} 결과는 실제로 받은 돈으로만 잰다. 예측이 틀렸다면 실제 IRR이 낮게 나와 그대로 드러난다.
\item 부도확률만 보지 않고 금리를 함께 보는 이유: LendingClub은 위험한 대출에 이미 높은 금리를 매긴다. 우리가 찾는 것은 위험에 비해 금리를 더 받는 대출이다.
\item 같은 달 안에서는 국채 금리가 모든 대출에 같으므로, 기준 수익률을 국채 대신 회사채로 바꿔도 줄 세우는 순서는 같다.
\end{itemize}

\section{근거: 승인비율별 결과}
\begin{table}[H]\centering\small
\begin{tabular}{@{}rcccc@{}}\toprule
 & \multicolumn{2}{c}{선별 가치: 무작위보다 더 번 \%p} & \multicolumn{2}{c}{Sharpe: 제안 $-$ 같은 비율 무작위}\\\cmidrule(lr){2-3}\cmidrule(l){4-5}
승인비율 & 값 (90\% 구간) & 더 번 달 & 값 (90\% 구간) & 판단\\\midrule
10\% & \textbf{+2.11} (1.80, 2.45) & 98\% & $-0.07$ ($-0.59$, 0.14) & 구분 안 됨\\
20\% & +1.74 (1.52, 1.99) & 100\% & +0.08 ($-0.47$, 0.40) & 구분 안 됨\\
30\% & +1.46 (1.29, 1.62) & 100\% & +0.23 ($-0.28$, 0.58) & 구분 안 됨\\
40\% & +1.19 (1.06, 1.33) & 100\% & +0.31 ($-0.12$, 0.61) & 구분 안 됨\\
50\% & +0.99 (0.84, 1.12) & 100\% & +0.28 ($-0.04$, 0.55) & 구분 안 됨\\
60\% & +0.82 (0.68, 0.94) & 100\% & +0.32 (0.08, 0.54) & \textbf{높음}\\
70\% & +0.63 (0.51, 0.74) & 100\% & +0.32 (0.15, 0.52) & \textbf{높음}\\
80\% & +0.48 (0.39, 0.56) & 100\% & +0.29 (0.15, 0.46) & \textbf{높음}\\
90\% & +0.28 (0.22, 0.33) & 100\% & +0.21 (0.14, 0.36) & \textbf{높음}\\\midrule
평균 & \textbf{+1.08} & & & \\\bottomrule
\end{tabular}
\caption{발행월 컷, 국채 기준, 수수료 1\%. 선별 가치는 승인 1달러당 연환산 IRR의 차이, 무작위는 10회 평균, 구간은 6개월 블록 부트스트랩 2000회(42단계, \ft{results/42_final_proposal_check.csv}). 평균은 무작위 반복 수에 따라 1.07--1.09\%p로 조금 다르다(38단계 1.07).}
\end{table}
\begin{enumerate}[leftmargin=*]
\item \textbf{모든 승인비율에서 무작위보다 더 벌었다.} 9개 승인비율 모두 구간이 0보다 위이고, 거의 모든 달에 더 벌었다.
\item \textbf{승인비율이 작을수록 1달러당 이점이 크다.} 많이 살수록 결국 거의 다 사게 되므로 차이가 줄어든다(80\%: +0.48\%p). 반대로 포트폴리오 전체 초과수익은 많이 살수록 크다(10\%: 0.65, 80\%: 4.34\%p, 위 상자).
\item \textbf{승인비율이 60\% 이상이면 덜 흔들리기도 한다.} 50\% 이하에서는 Sharpe가 무작위와 구분되지 않는다.
\item \textbf{돈의 출처:} 같은 등급 안에서 기대수익 점수 상위가 하위보다 연 +2.38\%p(90\% 구간 2.08--2.73), 48개월 모두 양수였다(7장).
\item \textbf{특정 모델 덕이 아니다:} XGBoost도 선별 가치가 같았고(1.03 vs 1.04\%p), 정답을 섞은 가짜 모델은 0 근처였다(7장).
\end{enumerate}

\section{부도확률 낮은 순(이전 최종안)과의 관계}
\begin{table}[H]\centering\small
\begin{tabularx}{\linewidth}{@{}>{\raggedright\arraybackslash}p{50mm}cc@{}}\toprule
 & \textbf{기대수익 순(제안)} & 부도확률 낮은 순\\\midrule
무작위보다 더 버나(선별 가치) & 9개 승인비율 모두 & 40\% 이상에서만, 평균 +0.20\%p\\
무작위보다 덜 흔들리나(Sharpe) & 60\% 이상에서 & 9개 승인비율 모두\\
80\%에서 Sharpe / 평균 초과수익 & 3.01 / 4.34\%p & 3.06 / 4.30\%p\\\bottomrule
\end{tabularx}
\end{table}
두 전략은 서로 다른 목표에서 강하다. 기대수익 순은 \textbf{더 벌기}, 부도확률 낮은 순은 \textbf{덜 흔들리기}다. 부도확률 낮은 순은 ``흔들림을 가장 줄이고 싶은 투자자''를 위한 대안으로 보고한다. 승인비율이 클 때(80\%)는 두 전략의 결과가 거의 같다.

\section{어떻게 여기까지 왔나}
\begin{enumerate}[leftmargin=*]
\item \textbf{v1(처음 규칙서):} 주 가설이 기대수익 순위의 선별 가치(H2)였다.
\item \textbf{v2.2(39단계):} ``과제가 Sharpe 최대화를 요구한다''는 전제에서 포트폴리오 Sharpe를 주 지표로 두고, 1-SE 규칙(1위 Sharpe에서 그 SE를 뺀 선 이상 중 승인비율 최대)으로 \textbf{부도확률 낮은 순 80\%}를 골랐다. Sharpe 1위는 10\%(3.82)였고 80\%는 3.06이었다.
\item \textbf{검토에서 나온 약점:} 1-SE 기준선 안의 후보를 ``1위와 통계적으로 구분되지 않는다''고 설명한 것은 틀렸다(10\% $-$ 80\% 차이 0.76, 구간 0.45--1.60). 또 V2의 Sharpe는 40--90\%에서 거의 평평해 80\%가 특별하지 않았다(90\%는 기준선보다 0.001 낮아 빠졌다). ``왜 80\%냐''에 답하기 어려웠다.
\item \textbf{목적 확정(2026-11-17, 18):} Sharpe 최고 튜닝은 목적이 아니고, 과제 공지를 다시 확인한 결과 Sharpe는 필수 지표가 아니었다. 그래서 투자자에게 더 직접적인 ``같은 양을 살 때 더 버는가''를 주 판정으로 두고, 처음 규칙서의 H2로 돌아갔다(test 규칙 v2.3, test를 열기 전 커밋과 태그).
\end{enumerate}
바꾸는 근거는 과제 해석이지만, train 결과를 모두 본 뒤의 결정이라는 사실은 기록한다. 판정 정의는 v1 그대로라 새로 만든 기준이 아니다.

\section{엄격 비판과 대응}
\begin{table}[H]\centering\small
\begin{tabularx}{\linewidth}{@{}>{\raggedright\arraybackslash}p{50mm}L@{}}\toprule
비판 & 점검 결과와 대응\\\midrule
``+1\%p''는 작은 승인비율이 끌어올린 평균이다 & 맞다. 10\%에서 +2.11, 80\%에서 +0.48\%p다(1달러당). 전체 돈의 초과수익은 반대로 10\%가 가장 작다(0.65\%p). 평균 하나가 아니라 승인비율별 표를 보인다. 많이 사는 투자자의 이점은 연 0.3--0.6\%p다.\\
승인비율이 작을 때는 더 흔들릴 수 있다 & 50\% 이하에서 Sharpe는 무작위와 구분되지 않고, 10\%는 점추정이 조금 낮다(2.55 vs 2.63). 금리가 높은 대출을 사기 때문이다. ``승인비율이 작을 때 덜 흔들린다는 근거는 없다''고 쓴다.\\
흔들림만 보면 부도확률 순이 낫다 & 맞다. 맞교환이며 표로 그대로 보인다.\\
$p$의 크기가 틀리면 순서가 바뀐다 & 모델 $p$는 실제보다 낮다(2013년 6.1\% vs 12.2\%). $p$를 실제 수준으로 맞춰 다시 골라도 승인 대상의 88--92\%가 같았다(5장).\\
결과를 보고 전략을 바꾼 것 아닌가 & H2는 v1의 주 가설이었고 정의도 그대로다. 바꾼 근거는 과제 해석이다. train 결과를 본 뒤라는 사실은 기록했다.\\
Sharpe 약 3은 너무 크다 & 만기 보유로 가격 변동이 없고, 2013--16년에 경기침체가 없어 부도가 몰리지 않았다. 우리 전략끼리 비교하는 데만 쓴다.\\
불경기에는? & 직접 확인하지 못했다. 팀 추가 점검(역 스트레스, 저장소 반영 전)에서는 부도가 모든 대출에서 고르게 늘면 약 6배까지 이점이 유지되지만, 모델이 안전하다고 본 절반의 부도만 약 2배가 되면 이점이 사라졌다. 이 제안은 ``불경기에도 모델의 위험 순서가 유지된다''는 가정에 기댄다.\\
살 수 있는 물량 & 개인이 조각으로 살 수 있는 대출 비중이 2013년 72\%에서 2016년 27\%로 줄었다. 상위부터 원하는 만큼 살 수 있다고 가정했다.\\\bottomrule
\end{tabularx}
\end{table}

\section{기준 수익률을 바꾸면}
기준을 국채보다 높게 잡아도 평균 초과수익 1위는 네 경우 모두 기대수익 순이다(38단계). 회사채(BAA) 기준에서는 Sharpe도 기대수익 순이 높다(승인 30\%: 1.29, 부도확률 순 0.92, 무작위 0.52). 부도확률이 낮은 대출은 금리도 낮아서, 기준이 오르면 남는 폭이 작아 조금만 나빠도 기준 아래로 떨어지기 때문이다.

\section{불확실성은 어디까지 셌나}
\begin{itemize}[leftmargin=*]
\item \textbf{센 것:} 기간 변화. 48개월을 6개월 이동 블록으로 다시 뽑았다(2000회, 모든 전략 공통 인덱스, 짝지은 차이). 그 표집 방식 아래의 불확실성일 뿐, 새로운 경기침체에서도 유지된다는 보증이 아니다.
\item \textbf{줄였지만 따로 안 센 것:} 무작위 선택(10회 평균).
\item \textbf{안 센 것:} 모델 seed(3개 평균), 대출 단위 표본오차.
\item 모든 숫자는 train 성과이므로 낙관적으로 평가됐을 가능성이 있다.
\end{itemize}
